\chapter{Caractéristiques du carnet d'ordres}\label{ch:rs:order-book-features}

Trois lots sont au meilleur cours acheteur et quarante au meilleur cours vendeur, la fourchette vaut un tick. Dans quel sens le prochain mouvement du mid-price va-t-il
se produire~? Le simulateur intrajournalier du carnet répond à partir de trois heures de ses propres données~: lorsque la file acheteuse est aussi petite, le mouvement suivant
est à la baisse trois fois sur quatre. Le carnet d'ordres est l'information la plus immédiate qu'un marché publie, et toute société de trading
qui agit à des horizons de quelques secondes la lit. Ce chapitre définit les caractéristiques standard du carnet d'ordres (déséquilibre au meilleur niveau et
au-delà, \omterm{def:rs:order-book-features:ofi}{déséquilibre du flux d'ordres}, \omterm{def:rs:order-book-features:micro}{mid pondéré} et \omterm{def:rs:order-book-features:micro}{microprix}, \omterm{def:rs:order-book-features:depletion}{taux d'épuisement de la file} et \omterm{def:rs:order-book-features:depletion}{taux d'annulation}), les mesure sur
firm.tape, compare le déséquilibre à la prédiction d'un modèle de files d'attente, et construit un moteur de caractéristiques en flux continu dans trois
langages. Ses données sont trois heures de messages simulés, soit 239\,528 messages.

\section{Déséquilibre au meilleur niveau}

\begin{definition}[Déséquilibre des files]\label{def:rs:order-book-features:imbalance}
Le \emph{déséquilibre des files}\index{déséquilibre des files} au meilleur niveau est $I = (q^b - q^a)/(q^b + q^a)$, où $q^b$ et $q^a$ sont les
quantités affichées au meilleur cours acheteur et au meilleur cours vendeur~; il appartient à $[-1, 1]$ et est positif lorsque la file acheteuse est la plus grande.
\end{definition}

\begin{omfigure}
\begin{tikzpicture}[font=\footnotesize,x=0.075cm,y=0.62cm]
\draw[black!60] (0,-0.5) -- (0,5.5);
\foreach \y/\p/\q in {4/100.02/35,3/100.01/28,2/100.00/40} {
  \fill[omThm!40] (0,\y-0.35) rectangle (\q,\y+0.35);
  \node[anchor=east] at (-2,\y) {\p};
  \node[anchor=west,omThm] at (\q+1,\y) {\q};
}
\foreach \y/\p/\q in {1/99.99/3,0/99.98/22,-1/99.97/31} {
  \fill[omDef!45] (0,\y-0.35) rectangle (\q,\y+0.35);
  \node[anchor=east] at (-2,\y) {\p};
  \node[anchor=west,omDef] at (\q+1,\y) {\q};
}
\node[anchor=west,omThm] at (48,3.4) {l'ask (lots)};
\node[anchor=west,omDef] at (48,0.2) {le bid (lots)};
\draw[black!70,dashed] (-14,1.5) -- (46,1.5);
\node[anchor=west,black!70] at (48,1.5) {fourchette~: un tick};
\node[anchor=west,align=left] at (48,-1.0) {déséquilibre au meilleur niveau $I = (3 - 40)/43 = -0.86$};
\end{tikzpicture}

\omcaption{Le carnet autour du meilleur niveau~: trois lots au meilleur cours acheteur (99.99), quarante au meilleur cours vendeur (100.00). Le \omterm{def:rs:order-book-features:imbalance}{déséquilibre des files} au
meilleur niveau vaut $-0.86$~; sur trois niveaux de chaque côté, il vaut $(56 - 103)/159 = -0.30$, parce que les demandes plus profondes sont plus grosses. Schéma, non
données.}\label{fig:rs:order-book-features:book}
\end{omfigure}

Lorsque la fourchette vaut un tick, le mid ne bouge que lorsque l'une des deux meilleures files se vide (ou lorsqu'un nouveau niveau apparaît à l'intérieur d'une
fourchette plus large). Une petite file acheteuse est plus proche de se vider qu'une grande file vendeuse, si bien que le déséquilibre prédit le sens du mouvement suivant.
Le premier modèle de ce phénomène est une course.

\begin{proposition}[Déséquilibre comme course entre deux files]\label{prop:rs:order-book-features:race}
Modélisons les deux meilleures files, en lots, comme des processus de naissance et de mort indépendants de taux d'arrivée $\lambda$ et de taux de départ $\mu$
(livre~4, chapitre~8), partant de $q^b$ et $q^a$. La probabilité que le prochain mouvement du mid soit à la hausse est la probabilité que la file vendeuse
se vide la première, $\int_0^\infty(1 - F_{q^b}(t))\,dF_{q^a}(t)$, où $F_q$ est la loi du temps que met une file partant de $q$ pour
se vider. Elle croît avec $q^b$, décroît avec $q^a$, et vaut $\tfrac12$ lorsque $q^b = q^a$.
\end{proposition}

\begin{proof}
Le mouvement suivant est à la hausse lorsque le côté vendeur se vide avant le côté acheteur~; les files étant indépendantes, c'est la course énoncée.
Une file de départ plus grande met stochastiquement plus de temps à se vider, ce qui donne la monotonie~; la symétrie donne $\tfrac12$.
\end{proof}

Sur les trois heures du simulateur, 2.03 lots par seconde arrivent à chaque meilleure file et 2.13 la quittent (ajouts au meilleur prix~;
annulations et exécutions à ce niveau). La \Cref{fig:rs:order-book-features:updown} confronte la prédiction de la course aux données, classe par
classe du déséquilibre. La direction est juste~: la probabilité d'un mouvement à la hausse passe de 0.245 dans le dixième
des observations où le carnet est le plus chargé côté vendeur, à 0.736 dans le dixième où il est le plus chargé côté acheteur. Le modèle est bien trop confiant~: il prédit 0.017 et 0.979. La raison tient au
mécanisme que le simulateur reproduit, et que les marchés réels partagent~: les files ne sont pas des marches aléatoires indépendantes, car le
prix efficient se déplace et les fournisseurs de liquidité placés du mauvais côté annulent leurs ordres. Le déséquilibre en est une vue bruitée, et un modèle de files d'attente
qui l'ignore surestime ce que cette vue vaut. Gould et Bonart (2016) ajustent la même relation par régression logistique sur
dix actions du Nasdaq et trouvent une amélioration fortement significative par rapport à un modèle nul, importante pour les actions à grand tick.

\begin{omfigure}
\begin{tikzpicture}
\begin{axis}[width=10cm,height=5.4cm,xlabel={déséquilibre des files $I$ (centre de la classe)},ylabel={P(le prochain mouvement du mid est à la hausse)},
  tick label style={font=\small},label style={font=\small},xmin=-1,xmax=1,ymin=0,ymax=1,grid=major,grid style={gray!25},
  legend columns=1,legend style={font=\footnotesize,at={(0.5,-0.32)},anchor=north,draw=none,fill=none,
  /tikz/every even column/.append style={column sep=8pt}},table/col sep=comma]
\addplot[omDef,very thick,mark=*] table[x=centre,y=empirical]{figdata/research/08-order-book-features/updown.csv};
\addplot[omThm,thick,dashed,mark=square] table[x=centre,y=model]{figdata/research/08-order-book-features/updown.csv};
\addplot[black!50,thin,sharp plot] coordinates {(-1,0.5) (1,0.5)};
\legend{traces simulées,course de deux files de naissance et de mort}
\end{axis}
\end{tikzpicture}

\omcaption{Probabilité que le prochain mouvement du mid-price soit à la hausse, par dixième de \omterm{def:rs:order-book-features:imbalance}{déséquilibre des files} à une fourchette d'un tick~: de 0.245 à 0.736 sur trois
heures de \texttt{firm.tape}, contre 0.017 à 0.979 pour la course de deux files de naissance et de mort indépendantes, avec les taux mesurés sur les
mêmes données. Données~: \texttt{firm.tape}, graine 8~; \texttt{firm.queues}.}\label{fig:rs:order-book-features:updown}
\end{omfigure}

\section{Profondeur au-delà du meilleur niveau}

\begin{definition}[Déséquilibre de profondeur, profil de profondeur]\label{def:rs:order-book-features:depth}
Le \emph{déséquilibre de profondeur}\index{déséquilibre de profondeur} sur $L$ niveaux est $(\sum_{i \le L}q^b_i - \sum_{i \le L}q^a_i)/(\sum_{i\le L}q^b_i
+ \sum_{i \le L}q^a_i)$, où $q^b_i, q^a_i$ sont les quantités aux $i$-ièmes meilleurs prix acheteur et vendeur. Le \emph{profil de profondeur}%
\index{profil de profondeur} d'un côté est le vecteur de ses quantités selon la distance au meilleur niveau.
\end{definition}

La profondeur au-delà du meilleur niveau est l'endroit d'où viendra le prochain meilleur cours une fois qu'une file se sera vidée, et où les gros ordres montrent leur jeu, ou
le cachent. Son information est plus lente que celle du meilleur niveau et plus facile à feindre (les ordres éloignés du meilleur niveau sont peu coûteux à placer et à annuler, ce à quoi
revient le livre~9, chapitre~29). Un \omterm{def:rs:order-book-features:depth}{profil de profondeur} est un vecteur de caractéristiques plutôt qu'un nombre~: sa pente, sa concentration, la
distance au premier niveau qui porte plus d'une quantité donnée. Le moteur de ce chapitre calcule le déséquilibre sur les cinq
premiers niveaux~; le \omterm{def:rs:order-book-features:depth}{profil de profondeur} lui-même est l'entrée des modèles appris par machine du livre~12.

\section{Déséquilibre du flux d'ordres}

\begin{definition}[Déséquilibre du flux d'ordres]\label{def:rs:order-book-features:ofi}
Soient $b_n, q^b_n, a_n, q^a_n$ le meilleur cours acheteur, sa taille, le meilleur cours vendeur et sa taille après le $n$-ième événement du carnet. La contribution de l'événement $n$ au \emph{déséquilibre du flux
d'ordres}\index{déséquilibre du flux d'ordres} est
\[
e_n = \mathbf 1_{b_n \ge b_{n-1}}q^b_n - \mathbf 1_{b_n \le b_{n-1}}q^b_{n-1} - \mathbf 1_{a_n \le a_{n-1}}q^a_n + \mathbf 1_{a_n \ge
a_{n-1}}q^a_{n-1},
\]
et le déséquilibre du flux d'ordres (OFI) sur un intervalle est la somme des $e_n$ qui s'y trouvent (Cont, Kukanov et Stoikov, 2014).
\end{definition}

Les quatre termes comptent la demande ajoutée au cours acheteur, la demande retirée de celui-ci, l'offre ajoutée au cours vendeur et l'offre retirée de celui-ci~; lorsqu'un
niveau de prix change, la file ancienne entière compte comme retirée et la nouvelle comme ajoutée. Cont, Kukanov et Stoikov ont constaté sur 50 actions américaines
que, sur de courts intervalles, les variations de prix sont principalement dictées par l'OFI, au moyen d'une relation linéaire dont la pente est inversement proportionnelle
à la profondeur du marché.

\begin{method}[Mesurer l'impact du flux d'ordres]\label{met:rs:order-book-features:ofi}
Sur une grille d'intervalles de longueur $\Delta$, régresser la variation du prix milieu (en ticks) sur l'OFI de l'intervalle divisé par la taille moyenne des meilleures files~;
la pente est l'impact sur le prix par unité de flux normalisé par la profondeur, et le $R^2$ indique la part de la variation du prix que le
flux explique à cet horizon.
\end{method}

Sur la bande simulée, la régression explique 10~\% des variations du prix milieu à une seconde, 53~\% à dix secondes et 79~\% à soixante~; la pente à dix
secondes vaut 0.47 tick par taille moyenne de la meilleure file de flux net (\cref{fig:rs:order-book-features:ofi}). À une seconde, la plupart des intervalles
ne présentent aucune variation du prix milieu, et ceux qui en présentent sont des variations d'un seul tick déclenchées par le dernier événement, que l'OFI ne perçoit qu'en partie.

\begin{omfigure}
\begin{tikzpicture}
\begin{axis}[width=9.5cm,height=5.2cm,xlabel={OFI sur 10 secondes (en tailles moyennes des meilleures files)},ylabel={variation du prix milieu (ticks)},
  tick label style={font=\small},label style={font=\small},grid=major,grid style={gray!25},table/col sep=comma]
\addplot[omDef,only marks,mark=*,mark size=1.8pt] table[x=ofi,y=dmid]{figdata/research/08-order-book-features/ofi.csv};
\addplot[omThm,thick,domain=-3.5:3.5,samples=2] {0.4664*x};
\end{axis}
\end{tikzpicture}

\omcaption{Variation du prix milieu sur des intervalles de dix secondes en fonction du \omterm{def:rs:order-book-features:ofi}{déséquilibre du flux d'ordres} de l'intervalle, en vingt classes d'effectifs égaux, avec
la droite de régression (pente 0.47 tick par unité, $R^2 = 0.53$ sur les 1\,078 intervalles). Données~: \texttt{firm.tape} à travers
\texttt{firm.lobfeat}, graine 8.}\label{fig:rs:order-book-features:ofi}
\end{omfigure}

\section{Le microprix}

\begin{definition}[Mid pondéré, microprix]\label{def:rs:order-book-features:micro}
Le \emph{mid pondéré}\index{mid pondéré} vaut $(a\,q^b + b\,q^a)/(q^b + q^a)$~: le milieu se déplace vers le côté de la file la plus petite. Le
\emph{microprix}\index{microprix} du carnet est le milieu augmenté de la taille et du signe espérés de la prochaine variation du prix milieu conditionnellement à
la classe de déséquilibre courante (à un écart d'un tick), estimés sur les données passées.
\end{definition}

Le \omterm{def:rs:order-book-features:micro}{mid pondéré} utilise le déséquilibre au moyen d'une formule fixe~; le \omterm{def:rs:order-book-features:micro}{microprix} apprend l'ajustement à partir des données. Stoikov (2018) construit
l'estimateur complet à partir de toute la suite des variations ultérieures du prix milieu et constate qu'il est un meilleur prédicteur des prix à court terme que le milieu ou le \omterm{def:rs:order-book-features:micro}{mid pondéré}~;
la version du livre ne conserve que le premier terme de cette suite. Ajusté sur la première moitié de la session simulée et évalué sur la seconde, il présente une erreur quadratique moyenne
inférieure de 10~\% à celle du prix milieu pour la prévision à une seconde~; le \omterm{def:rs:order-book-features:micro}{mid pondéré} est moins bon de 27~\% à une seconde, mais le meilleur à cinq
et trente secondes (inférieur de 12~\% et de 3~\% au prix milieu). Une correction d'ordre un est juste pour le mouvement suivant et aveugle au-delà~;
le \omterm{def:rs:order-book-features:micro}{mid pondéré} est grossier, mais il continue d'indiquer la bonne direction.

\begin{omfigure}
\begin{tikzpicture}
\begin{axis}[width=8.5cm,height=5cm,ybar,bar width=14pt,ylabel={EQM relative au prix milieu},tick label style={font=\small},
  label style={font=\small},xmin=-0.6,xmax=2.6,ymin=0,ymax=1.4,xtick={0,1,2},xticklabels={1 seconde,5 secondes,30 secondes},
  grid=major,grid style={gray!25},legend columns=2,legend style={font=\footnotesize,at={(0.5,-0.25)},anchor=north,draw=none,
  fill=none,/tikz/every even column/.append style={column sep=8pt}},table/col sep=comma]
\addplot[fill=omThm!45,draw=omThm] table[x=k,y=wmid]{figdata/research/08-order-book-features/forecast.csv};
\addplot[fill=omDef!55,draw=omDef] table[x=k,y=micro]{figdata/research/08-order-book-features/forecast.csv};
\addplot[black!70,thin,sharp plot,forget plot] coordinates {(-0.6,1) (2.6,1)};
\legend{mid pondéré,microprix (un pas)}
\end{axis}
\end{tikzpicture}

\omcaption{Erreur quadratique moyenne du \omterm{def:rs:order-book-features:micro}{mid pondéré} et du \omterm{def:rs:order-book-features:micro}{microprix} à un pas, comme prévisions du prix milieu 1, 5 et 30 secondes plus tard,
relativement à l'erreur propre du prix milieu, sur la seconde moitié de la session simulée (la table du \omterm{def:rs:order-book-features:micro}{microprix} est ajustée sur la première). Données~:
\texttt{firm.tape} à travers \texttt{firm.lobfeat}, graine 8.}\label{fig:rs:order-book-features:forecast}
\end{omfigure}

\section{Épuisement des files et annulations}

\begin{definition}[Taux d'épuisement de la file, taux d'annulation]\label{def:rs:order-book-features:depletion}
Le \emph{taux d'épuisement de la file}\index{taux d'épuisement de la file} d'une meilleure file est la quantité qui la quitte par seconde (annulations et
exécutions) sur une fenêtre glissante~; son \emph{taux d'annulation}\index{taux d'annulation} ne compte que les annulations.
\end{definition}

Les taux portent ce que les tailles ne portent pas~: une file de quarante lots qui perd dix lots par seconde est plus fragile qu'une file de vingt lots qui n'en perd aucun. Dans le simulateur,
les fournisseurs de liquidité annulent plus vite du côté que le prix efficient a quitté, comme le font les véritables lorsque leurs cotations deviennent périmées.
La différence entre les annulations au meilleur cours vendeur et celles au meilleur cours acheteur sur la dernière seconde constitue donc une vue directe du prix caché~:
sa corrélation avec le sens de la prochaine variation du prix milieu vaut 0.30, contre 0.25 pour le \omterm{def:rs:order-book-features:imbalance}{déséquilibre des files}. La leçon générale est celle du chapitre~:
les caractéristiques les plus informatives sont celles qui sont les plus proches du mécanisme qui fait bouger le prix.

\section{Fiches de prédicteur}

\begin{predictorcard}{Déséquilibre des files}\label{pred:rs:order-book-features:imbalance}
\sfield{Définition} $I = (q^b - q^a)/(q^b + q^a)$ au premier niveau, après chaque message.
\sfield{Entrées et horodatages} Le carnet après le dernier message reçu~; l'heure de réception est le \omterm{def:rs:point-in-time-data-and-the-biases:bitemporal}{temps de connaissance}.
\sfield{Justification} La plus petite file se vide en premier (la \cref{prop:rs:order-book-features:race})~; les fournisseurs de liquidité situés du mauvais côté du
prix efficient se retirent.
\sfield{Horizon et demi-vie} La prochaine variation du prix milieu, en secondes. P(hausse) passe de 0.245 à 0.736 selon les dixièmes de $I$ sur la bande simulée.
\sfield{Normalisation} Aucune nécessaire~: bornée~; on la découpe en classes ou on la fournit à un modèle logistique.
\sfield{Modes de défaillance} Tailles factices (spoofing) (livre~9, chapitre~29)~; ordres cachés et icebergs~; fourchettes larges, où un mouvement n'a pas besoin de vider une file.
\sfield{Sources} Gould et Bonart (2016)~; \texttt{rs\_orderbook.up\_probability}.
\end{predictorcard}

\begin{predictorcard}{Déséquilibre du flux d'ordres (OFI)}\label{pred:rs:order-book-features:ofi}
\sfield{Définition} La somme des $e_n$ sur les $\Delta$ dernières secondes, divisée par la taille moyenne de la meilleure file.
\sfield{Entrées et horodatages} Chaque événement du carnet, dans l'ordre, avec les heures de réception.
\sfield{Justification} La demande nette au premier niveau fait bouger le prix linéairement, avec une pente inversement proportionnelle à la profondeur.
\sfield{Horizon et demi-vie} Simultané dans sa définition~; en tant que prédicteur, les secondes suivantes. $R^2$ de 10~\%, 53~\% et 79~\% à 1, 10 et 60
secondes (simultané) sur la bande simulée.
\sfield{Normalisation} Par la profondeur~; par la volatilité entre instruments.
\sfield{Modes de défaillance} Il faut chaque événement dans l'ordre (un trou dans le flux fausse la somme)~; les ajouts par spoofing comptent comme demande.
\sfield{Sources} Cont, Kukanov et Stoikov (2014)~; \texttt{rs\_orderbook.ofi\_regression}.
\end{predictorcard}

\begin{predictorcard}{Déséquilibre des annulations de cotations périmées}\label{pred:rs:order-book-features:cancel}
\sfield{Définition} Lots annulés au meilleur cours vendeur moins lots annulés au meilleur cours acheteur sur la dernière seconde.
\sfield{Entrées et horodatages} Messages d'annulation avec leurs prix et heures de réception.
\sfield{Justification} Les fournisseurs de liquidité retirent les cotations que le prix efficient a laissées derrière lui.
\sfield{Horizon et demi-vie} La prochaine variation du prix milieu~; corrélation de 0.30 avec son sens sur la bande simulée.
\sfield{Normalisation} Par la taille des files, ou classé entre instruments.
\sfield{Modes de défaillance} Annulation et remplacement au même prix par un algorithme qui rafraîchit ses cotations~; différences entre plateformes
dans la manière dont les modifications sont déclarées.
\sfield{Sources} Le mécanisme du simulateur (chapitre~2)~; \texttt{rs\_orderbook.cancel\_feature}.
\end{predictorcard}

\section{Tutoriel~: lire le livre}\label{tut:rs:order-book-features:read}

\begin{tutorial}
\textbf{Objectif.} Faire passer trois heures de messages simulés dans le moteur de caractéristiques, mesurer le déséquilibre, l'OFI et le \omterm{def:rs:order-book-features:micro}{microprix}, et comparer le
déséquilibre avec le modèle de files d'attente. \textbf{État final~:}
les \cref{fig:rs:order-book-features:updown,fig:rs:order-book-features:ofi,fig:rs:order-book-features:forecast}.
\begin{tutsteps}
\item \textbf{Le moteur.} Appliquer le message aux tables d'ordres et de niveaux, lire le nouveau sommet du carnet, calculer la contribution
de l'OFI à partir du meilleur niveau précédent, puis les déséquilibres, le \omterm{def:rs:order-book-features:micro}{mid pondéré} et le \omterm{def:rs:order-book-features:micro}{microprix}.
\omcode{code/firm/lobfeat/firm_lobfeat.py}{66}{96}{Un message traité par le moteur de référence Python.}\label{lst:rs:order-book-features:engine}
\item \textbf{La course.} La probabilité que la file du meilleur cours vendeur se vide en premier, selon les lois d'épuisement du livre~4 et de
\texttt{firm.queues}, conditionnée au vidage d'une des deux files dans les dix minutes.
\omcode{code/research/08-order-book-features/python/rs_orderbook.py}{79}{86}{La course du processus de naissance et de mort.}\label{lst:rs:order-book-features:race}
\item \textbf{Exécuter} \texttt{up\_probability()}, \texttt{best\_level\_rates()}, \texttt{ofi\_regression} pour 1, 10 et 60 secondes,
\texttt{forecast\_errors()}, \texttt{cancel\_feature()} et \texttt{fig\_orderbook.py}.
\end{tutsteps}
\textbf{Pour aller plus loin.} Ajuster une régression logistique du sens du prochain mouvement sur $I$ et sur le déséquilibre des annulations conjointement~;
estimer aussi la table du \omterm{def:rs:order-book-features:micro}{microprix} pour des fourchettes de deux ticks.
\end{tutorial}

\section{Construction~: le moteur de caractéristiques en flux continu}\label{bld:rs:order-book-features:lobfeat}

\begin{build}
\bfield{Rôle} Les caractéristiques du carnet d'ordres de la société miniature, calculées message par message de la même façon dans la recherche (Python) et dans
la production (C++20, avec un jumeau en Rust)~; le teneur de marché du livre~11 et le modèle de carnet d'ordres du livre~12 les utilisent.
\bfield{Interface} \texttt{Engine(levels, g).on(kind, oid, side, price, qty)} renvoyant \texttt{Features} (meilleurs cours et tailles,
\texttt{imbalance}, \texttt{depth\_imbalance}, \texttt{ofi}, \texttt{ofi\_cum}, \texttt{wmid}, \texttt{micro})~; \texttt{run(msgs)}~;
\texttt{LobFeatures} dans \texttt{cpp/firm\_lobfeat.hpp} et \texttt{rust/src/lib.rs}.
\bfield{Règles} Prix et tailles entiers~; aucune caractéristique tant que chaque côté n'a pas de cotation~; l'OFI du premier événement à deux côtés est nul~; arithmétique identique
dans les trois langages.
\bfield{Tests d'acceptation} \texttt{code/firm/lobfeat/tests/}~: une suite de messages d'ajout, d'exécution et de changement de niveau vérifiée à la main~; le jeu de test versionné
reproduit par la référence Python~; déséquilibres bornés et OFI suivant le prix milieu~; les tests C++ et Rust rejouent le jeu de 2\,769 messages
et retrouvent chaque caractéristique des 2\,766 lignes à deux côtés.
\bfield{Pour aller plus loin} \omterm{def:rs:order-book-features:depth}{Profils de profondeur} sous forme de vecteurs de taille fixe~; taux d'épuisement et \omterm{def:rs:order-book-features:depletion}{taux d'annulation} sur fenêtre glissante dans le moteur~; une file
à écriture unique et sans verrou qui l'alimente (livre~13).
\end{build}

Le cœur C++, la seule partie délicate du moteur, est la mise à jour du flux d'ordres à partir du meilleur niveau du carnet précédent~:

\omcode{code/firm/lobfeat/cpp/firm_lobfeat.hpp}{43}{60}{Le moteur C++20 après chaque message.}\label{lst:rs:order-book-features:cpp}

\begin{omsources}
\item R.~Cont, A.~Kukanov et S.~Stoikov, «~The price impact of order book events~», \emph{Journal of Financial Econometrics} 12(1), 2014.
\item S.~Stoikov, «~The micro-price: a high-frequency estimator of future prices~», \emph{Quantitative Finance} 18(12), 2018.
\item M.~D.~Gould et J.~Bonart, «~Queue imbalance as a one-tick-ahead price predictor in a limit order book~», \emph{Market Microstructure
and Liquidity} 2(2), 2016.
\item R.~Cont, S.~Stoikov et R.~Talreja, «~A stochastic model for order book dynamics~», \emph{Operations Research} 58(3), 2010.
\item W.~Huang, C.-A.~Lehalle et M.~Rosenbaum, «~Simulating and analyzing order book data: the queue-reactive model~», \emph{Journal of the
American Statistical Association} 110(509), 2015.
\end{omsources}

\section{Exercices}

\begin{exercise}[$\star$]\label{exo:rs:order-book-features:1}
Calculer le \omterm{def:rs:order-book-features:imbalance}{déséquilibre des files}, le prix milieu et le \omterm{def:rs:order-book-features:micro}{mid pondéré}, sachant que le meilleur cours acheteur est de 99.98 pour 700 actions et le meilleur cours vendeur de 99.99 pour 300.
\end{exercise}

\begin{exercise}[$\star$]\label{exo:rs:order-book-features:2}
Le cours acheteur passe de 99.98 (500 actions) à 99.99 (200 actions) et le cours vendeur reste à 100.00 (400 actions, inchangé).
Calculer la contribution de l'événement à l'OFI.
\end{exercise}

\begin{exercise}[$\star$]\label{exo:rs:order-book-features:3}
La pente de l'OFI vaut 0.47 tick par taille moyenne de la meilleure file, cette taille étant de 2\,750 actions. Calculer le mouvement du prix milieu prédit par un flux acheteur net de 11\,000
actions sur dix secondes.
\end{exercise}

\begin{exercise}[$\star\star$]\label{exo:rs:order-book-features:4}
Pour la course avec $\lambda = \mu$ et des files d'un lot au cours acheteur et $n$ lots au cours vendeur, montrer que la probabilité d'un mouvement haussier tend vers
zéro lorsque $n$ croît. Que dit cela des tranches extrêmes du modèle~?
\end{exercise}

\begin{exercise}[$\star\star$]\label{exo:rs:order-book-features:5}
Pourquoi le $R^2$ de l'OFI augmente-t-il avec la longueur de l'intervalle, de 10~\% pour une seconde à 79~\% pour soixante secondes~?
\end{exercise}

\begin{exercise}[$\star\star$]\label{exo:rs:order-book-features:6}
Montrer que le \omterm{def:rs:order-book-features:micro}{mid pondéré} égale le prix milieu plus $\tfrac s2 I$, avec $s$ la fourchette et $I$ le \omterm{def:rs:order-book-features:imbalance}{déséquilibre des files}.
\end{exercise}

\begin{exercise}[$\star\star\star$]\label{exo:rs:order-book-features:7}
\emph{Programmation.} En utilisant la bande du chapitre, calculer la corrélation du \omterm{def:rs:order-book-features:depth}{déséquilibre de profondeur} sur cinq niveaux avec le sens du
prochain mouvement du prix milieu, et la comparer à celle du déséquilibre au meilleur niveau.
\end{exercise}

\begin{exercise}[$\star\star\star$]\label{exo:rs:order-book-features:8}
\emph{Trouver la faille.} «~Notre modèle de carnet d'ordres prédit le prochain mouvement du prix milieu avec une justesse de 74~\%. Nous avons calculé les caractéristiques à partir du carnet consolidé
horodaté avec les horodatages de la bourse et la cible tirée du même flux.~»
\end{exercise}

\section{Problème~: trois lots contre quarante}

\begin{problem}[{Problème du week-end --- combien le carnet d'ordres sait-il, et combien un modèle de files croit-il savoir~?}]\label{pb:rs:order-book-features:1}
Trois heures de \texttt{firm.tape} (graine 8), 239\,528 messages, passés par \texttt{firm.lobfeat}.

\medskip\noindent\textbf{Partie I --- Déséquilibre.}
\begin{enumerate}
\item Quelles sont les probabilités d'un mouvement haussier dans les déciles inférieur et supérieur du \omterm{def:rs:order-book-features:imbalance}{déséquilibre des files}~?
\item Combien d'observations tombent dans chacun de ces deux déciles~?
\item Quelles sont les tailles moyennes des files (lots) dans les déciles extrêmes~?
\item Pour quel déséquilibre la probabilité vaut-elle un demi~?
\end{enumerate}

\medskip\noindent\textbf{Partie II --- Le modèle.}
\begin{enumerate}[resume]
\item Quels taux d'arrivée et de départ (lots par seconde) possède chaque meilleure file~?
\item Que prédit la course pour les déciles extrêmes~?
\item Pourquoi le modèle est-il trop confiant~?
\item Quel mécanisme du simulateur un modèle correct devrait-il inclure~?
\end{enumerate}

\medskip\noindent\textbf{Partie III --- Flux.}
\begin{enumerate}[resume]
\item Quels sont les $R^2$ de l'OFI à 1, 10 et 60 secondes, et sa pente à dix secondes~?
\item Quelle est la taille moyenne de la meilleure file utilisée pour la normaliser~?
\item Comment le \omterm{def:rs:order-book-features:micro}{microprix} à un pas se compare-t-il au mid à une seconde~?
\item Quelle prévision l'emporte à cinq et à trente secondes, et de combien~?
\item Quelle est la corrélation du déséquilibre des annulations avec le mouvement suivant, et celle du \omterm{def:rs:order-book-features:imbalance}{déséquilibre des files}~?
\end{enumerate}

\medskip\noindent\textbf{Partie IV --- Ingénierie et jugement.}
\begin{enumerate}[resume]
\item Pourquoi les trois implémentations du moteur doivent-elles produire des résultats identiques~?
\item Combien de lignes de référence les tests C++ et Rust vérifient-ils~?
\item Quelles caractéristiques exigent chaque message dans l'ordre, et lesquelles survivent à un flux d'instantanés~?
\item Laquelle des caractéristiques du chapitre se laisserait déplacer le plus facilement par un spoofeur~?
\item Énoncer le \emph{résultat nommé}~: les probabilités empirique et modélisée d'un mouvement à la hausse dans les déciles extrêmes de déséquilibre, et le
$R^2$ de l'OFI à dix secondes.
\item Qu'ajouteriez-vous en premier au modèle de course~?
\item En une phrase~: que vous dit le carnet d'ordres~?
\end{enumerate}
\end{problem}

\section{Questions d'entretien}

\begin{interviewq}[$\star$ \iqroles{trader, chercheur}]\label{iq:rs:order-book-features:1}
La file acheteuse est dix fois plus longue que la file vendeuse à une fourchette d'un tick. Que prévoyez-vous pour le prochain mouvement de prix, et pourquoi~?
\end{interviewq}

\begin{interviewq}[$\star\star$ \iqroles{chercheur}]\label{iq:rs:order-book-features:2}
Définissez le \omterm{def:rs:order-book-features:ofi}{déséquilibre du flux d'ordres} et expliquez pourquoi il est relié linéairement aux variations de prix.
\end{interviewq}

\begin{interviewq}[$\star\star$ \iqroles{chercheur, développeur}]\label{iq:rs:order-book-features:3}
Qu'est-ce qu'un \omterm{def:rs:order-book-features:micro}{microprix}~? Comment en estimeriez-vous un~?
\end{interviewq}

\begin{interviewq}[$\star\star$ \iqroles{développeur}]\label{iq:rs:order-book-features:4}
Vos caractéristiques de recherche sont en Python et la production est en C++. Comment garantissez-vous qu'elles calculent la même chose~?
\end{interviewq}

\begin{interviewq}[$\star\star$ \iqroles{trader, chercheur}]\label{iq:rs:order-book-features:5}
Comment une caractéristique du carnet d'ordres peut-elle être manipulée, et comment rendriez-vous un modèle robuste à cette manipulation~?
\end{interviewq}

\begin{interviewq}[$\star\star\star$ \iqroles{chercheur}]\label{iq:rs:order-book-features:6}
Modélisez les files du meilleur cours acheteur et du meilleur cours vendeur comme des processus de naissance et de mort indépendants. Que prédit le modèle pour le prochain mouvement, et que
laisse-t-il de côté~?
\end{interviewq}
